%%%%%%%%%%%%%%%%%%%%%%%%%%%%%%%%%%%%%%%%%%%%%%%%%%%%%%%%%%%%%%%%%%%%%%
% Lớp 10 | Mệnh đề. Tập hợp | Các phép toán tập hợp
% Dạng 0D1?3-4: Hiệu và phần bù (dùng đoạn, khoảng)
% Số câu: 26  (sinh tự động bởi tach_cau.py; từ nay thầy sửa trực tiếp file này)
%%%%%%%%%%%%%%%%%%%%%%%%%%%%%%%%%%%%%%%%%%%%%%%%%%%%%%%%%%%%%%%%%%%%%%

%%%%%%%%%%%%%%%% TRẮC NGHIỆM 4 PHƯƠNG ÁN %%%%%%%%%%%%%%%%

% nguồn: MaHoaNH_TracNghiem_Lop10
\begin{ex}%[0D1N3-4]%[ma:59bd772c]
	Cho hai tập hợp $A=(1; 5]$,  $B=(2; 7]$. Tập hợp $A \setminus B$ là
	\choice 
		{\True $(1; 2]$} 
		{$(2; 5)$} 
		 {$(-1; 7]$}
		 {$(-1; 2)$}
	\loigiai{ 
		Ta có $A  \setminus B=(1; 2]$.
	 }
\end{ex}

% nguồn: MaHoaNH_TracNghiem_Lop10
\begin{ex}%[0D1H3-4]%[ma:56f74312]
	Cho hai tập hợp $A=\{1;2;5;6\}$, $B=\{1;2;3;4;5;6;7;8\}$ khi đó tập $\mathrm{C}_{B}A$ là 
	\choice  {$\{1;2;4;6\}$}{$\{4;6\}$}{\True$\{3;4;7;8\}$}{$\{2;6;7;8\}$}
	\loigiai{ 
		Tập hợp $\mathrm{C}_{B}A$ (phần bù của $A$ trong $B$) là tập hợp tất cả  các phần tử có trong tập $B$ mà tập $A$ không có.
		$\mathrm{C}_{B}A = B \setminus A  = \{1;2;3;4;5;6;7;8\} \setminus \{1;2;5;6\} = \{3;4;7;8\}$.
	 }
\end{ex}

% nguồn: MaHoaNH_TracNghiem_Lop10
\begin{ex}%[0D1H3-4]%[ma:de61549b]
	Cho tập hợp $A=\{x\in\mathbb{N}|3x-5>10\}$ khi  đó
	\choice  {$\mathrm{C}_{\mathbb{N}}A=\{1;2;3;4;5\}$}{\True $\mathrm{C}_{\mathbb{N}}A=\{0;1;2;3;4;5\}$}{$\mathrm{C}_{\mathbb{N}}A=\{1;2;3\}$}{$\mathrm{C}_{\mathbb{N}}A=\{1;2;5\}$}
	\loigiai{ 
		Giải bất  phương trình $3x-5>10\Leftrightarrow x>5$.\\
		Vì $x\in\mathbb{N}$ nên  $A=\{6;7;8;9;10;\ldots\}$.\\
		Khi  đó $\mathrm{C}_{\mathbb{N}}A=\mathbb{N}\setminus A=\{0;1;2;3;4;5\}$.
	} 
\end{ex}

% nguồn: MaHoaNH_TracNghiem_Lop10
\begin{ex}%[0D1H3-4]%[ma:d3c7dcd9]
	Cho ba tập hợp $A=[-2; 2]$, $B=[1; 5]$, $C=[0; 1)$.  Khi đó tập $(A \setminus B) \cap C$ là
	 \choice
		{$\{0;  1\}$}
		{\True $[0;  1)$}
		{$(-2;1)$} 
		 {$[-2; 5]$}
	\loigiai{ 
		Ta có $A \setminus B=[-2; 1) \Rightarrow(A \setminus  B) \cap C=[0; 1)$.
	 }
\end{ex}

% nguồn: MaHoaNH_TracNghiem_Lop10
\begin{ex}%[0D1H3-4]%[ma:ece9acf4]
	Cho hai tập hợp $M=\left\{x\in \mathbb{R}\mid -1\le x\le 4\right\}$ và $\mathrm{C}_{\mathbb{R}}  N=\left(-\infty;0\right)$. Hãy tìm tập hợp $M\setminus N$.
	 \choice
		 {$\left[-1;4\right]$}
		 {\True $\left[-1;0\right)$}
		 {$\left[0;4\right]$}
		 {$\left(4;+\infty \right)$}
	 \loigiai{
		Ta có $M=\left[-1;4\right]$; $N=\left[0;+\infty \right)$ suy ra $M\setminus  N=\left[-1;0\right)$.
	} 
\end{ex}

% nguồn: MaHoaNH_TracNghiem_Lop10
\begin{ex}%[0D1H3-4]%[ma:a0c8c823]
	Cho các tập hợp $A=\{x \in \mathbb{R} \mid  x<3\}$, $B=\{x \in \mathbb{R} \mid 1<x \leq 5\}$, \\$C=\{x \in \mathbb{R} \mid-2 \leq x \leq 4\}$. Khi đó $(B \cup C) \setminus(A \cap C)$ bằng
	\choice 
		 {$[-2; 3)$}
		{\True $[3; 5]$} 
		 {$(-\infty; 1]$}
		{$[-2;  5]$}
	\loigiai{ 
		Ta có $A=(-\infty; 3)$,  $B=(1; 5]$, $C=[-2; 4]$.\\
		$(B \cup C) \setminus(A \cap C)=\{(1; 5] \cup[-2; 4]\} \setminus\{(-\infty; 3) \cap[-2; 4]\}=[-2;  5] \setminus[-2; 3)=[3; 5]$.
	 }
\end{ex}

% nguồn: MaHoaNH_TracNghiem_Lop10
\begin{ex}%[0D1H3-4]%[ma:5c4579d5]
	Cho các tập $A=\{x \in \mathbb{R} \mid x \geq-1\}$, $B=\{x \in \mathbb{R} \mid x<3\}$. Tập $\mathbb{R} \setminus(A \cap  B)$ là
	 \choice
		 {\True $(-\infty;-1) \cup[3;+\infty)$}
		{$(-1;  3]$}
		 {$[-1; 3)$}
		{$(-\infty;-1] \cup(3;+\infty)$} 
	 \loigiai{
		Ta có  $A=[-1;+\infty)$, $B=(-\infty; 3)$.\\
		Khi đó $A \cap B=[-1; 3)  \Rightarrow \mathbb{R} \setminus(A \cap B)=(-\infty;-1) \cup[3;+\infty)$.
	 }
\end{ex}

% nguồn: MaHoaNH_TracNghiem_Lop10
\begin{ex}%[0D1H3-4]%[ma:0578f957]
	Cho hai tập hợp $A=\left(\sqrt{2};+\infty\right)$ và $B=\left(-\infty;  \dfrac{\sqrt{5}}{2}\right]$. Khi đó $(A \cap B) \cup(B \setminus A)$ là
	 \choice
		{$\left[\dfrac{\sqrt{5}}{2};  \sqrt{2}\right]$}
		 {$(\sqrt{2};+\infty)$}
		{\True $\left(-\infty; \dfrac{\sqrt{5}}{2}\right]$} 
		 {$\left(-\infty; \dfrac{\sqrt{5}}{2}\right)$}
	 \loigiai{
		Ta có $A \cap B=\varnothing$, $B \setminus A=\left(-\infty; \dfrac{\sqrt{5}}{2}\right]$.\\ 
		Do đó $(A \cap B) \cup(B \setminus  A)=\left(-\infty; \dfrac{\sqrt{5}}{2}\right]$.
	} 
\end{ex}

% nguồn: MaHoaNH_TracNghiem_Lop10
\begin{ex}%[0D1H3-4]%[ma:564a6777]
	Cho tập hợp $C_{\mathbb{R}} A=\left[-3; \sqrt{8}\right)$,  $C_{\mathbb{R}} B=(-5; 2) \cup \left(\sqrt{3}; \sqrt{11}\right)$. Tập $C_{\mathbb{R}}(A \cap B)$ là
	\choice 
		{$\left(-3; \sqrt{3}\right)$} 
		 {$\varnothing$}
		 {\True $\left(-5; \sqrt{11}\right)$}
		{$(-3; 2) \cup \left(\sqrt{3}; \sqrt{8}\right)$} 
	\loigiai{ 
		Ta có $C_{\mathbb{R}} A=\left[-3; \sqrt{8}\right)$, $C_{\mathbb{R}} B=(-5; 2) \cup  \left(\sqrt{3}; \sqrt{11}\right)=\left(-5; \sqrt{11}\right)$.\\
		$A=(-\infty;-3) \cup  \left[\sqrt{8};+\infty\right)$, $B=(-\infty;-5] \cup \left[\sqrt{11};+\infty\right)$.\\
		Suy ra  $A \cap B=(-\infty;-5] \cup \left[\sqrt{11};+\infty\right)$.\\
		Vậy $C_{\mathbb{R}}(A  \cap B)=\left(-5; \sqrt{11}\right)$.
	 }
\end{ex}

% nguồn: MaHoaNH_TracNghiem_Lop10
\begin{ex}%[0D1H3-4]%[ma:437037f3]
	Cho tập hợp  $A=\{x \in \mathbb{R} \mid 2 \leqslant x<5\}$. Phần bù của tập hợp $A$ trong $\mathbb{R}$ là tập nào sau đây?
	 \choice
		 {$[5;+\infty)$}
		{$(-\infty;  2)$}
		{$(-\infty; 2] \cup(5;+\infty)$} 
		{\True $(-\infty;  2) \cup[5;+\infty)$}
	\loigiai{ 
		Ta  có $A = \{x \in \mathbb{R} \mid 2 \leqslant x<5\}
		= [2;5)$.\\ 
		Khi đó phần bù của  tập hợp $A$ trong $\mathbb{R}$ là
		 $\mathrm{C_{\mathbb{R}}A}
		= \mathbb{R} \setminus A = (-\infty; 2) \cup[5;+\infty)$. 
	} 
\end{ex}

% nguồn: MaHoaNH_TracNghiem_Lop10
\begin{ex}%[0D1Y3-4]%[ma:3113212d]
	Một cửa hàng bán áo sơ mi thống kê  số lượng áo bán ra trong tháng $ 1 $ như bảng sau.
	\begin{center}
		\begin{tabular}{|c|l|l|l|l|l|l|}
			\hline Cơ áo & 37 & 38 & 39 & 40 & 41 &  42 \\
			\hline Số lượng & 35 & 42 & 50  & 38 & 32 & 48 \\
			\hline 
		\end{tabular}
	\end{center}
	Mốt của mẫu số liệu trên  bằng?
	\choice 
		 {$ 42 $}
		{$ 39 $} 
		 {$ 41 $}
		{\True $  50 $}
	 \loigiai{
		Mốt của mẫu số  liệu trên là $ 50 $.
	 }
\end{ex}


%%%%%%%%%%%%%%%% ĐÚNG/SAI %%%%%%%%%%%%%%%%

% nguồn: MaHoaNH_DungSai_Lop10
\begin{ex}%[0D1N3-4]%[ma:d60d7622]
	Cho hai tập hợp $A=(-1;+\infty)$,  $B=(-\infty;-1]$.
	 \choiceTF
		 {\True $A\backslash B=(-1;+\infty)$}
		{\True $B\backslash A=(-\infty;-1]$} 
		{  $C_\mathbb{R}A=(-\infty;-1)$}
		 {\True $C_\mathbb{R}B=(-1;+\infty)$}
	\loigiai{ 
		\begin{itemchoice}
			 \itemch $A\backslash B=(-1;+\infty)$.
			\itemch $B\backslash  A=(-\infty;-1]$.
			\itemch  $C_\mathbb{R}A=(-\infty;-1]$.
			\itemch $C_\mathbb{R}B=(-1;+\infty)$. 
		\end{itemchoice}
	 }
\end{ex}

% nguồn: MaHoaNH_DungSai_Lop10
\begin{ex}%[0D1N3-4]%[ma:a5213f07]
	Cho hai tập hợp $A=\{x\in\mathbb{R}||x|<3\}$,  $B=[-2;2]$.
	\choiceTF 
		{ $A\backslash B=\varnothing$} 
		 {\True $B\backslash A=\varnothing$}
		{\True $C_\mathbb{R}A=(-\infty;-3]\cup  [3;+\infty)$}
		{ $(C_\mathbb{R}B)\cap A=(-\infty;-2)\cup (2;+\infty)$} 
	 \loigiai{$A=(-3;3)$, $B=[-2;2]$.
		\begin{itemchoice}
			\itemch $A\backslash B=(-3;-2)\cup  (2;3)$.
			\itemch $B\backslash A=\varnothing$. 
			\itemch $C_\mathbb{R}A=(-\infty;-3]\cup  [3;+\infty)$.
			 \itemch $(C_\mathbb{R}B)\cap A=(-3;-2)\cup (2;3)$.
		\end{itemchoice}
	} 
\end{ex}

% nguồn: MaHoaNH_DungSai_Lop10
\begin{ex}%[0D1H3-4]%[ma:434ed489]
	Cho hai tập hợp $A=(-\infty;3]$ và $B=[2;5)$. Xét tính đúng  sai của các mệnh đề sau:
	\choiceTF {\True $A\cap B=[2;3]$}{\True $A\cup B=(-\infty;5)$}{ Tập hợp $A\setminus B=(-\infty;2]$ có 3 số nguyên không âm}{ Tập hợp $C=\{x|x\in\mathbb{N},x\in B\setminus A\}$  có 4 tập con}
	\loigiai{ 
		\begin{itemchoice}
			\itemch $A\cap B$ là giao  của $(-\infty;3]$ và $[2;5)$, ta được $[2;3]$.
			\itemch $A\cup B$ là hợp của $(-\infty;3]$  và $[2;5)$, ta được $(-\infty;5)$.
			\itemch $A\setminus B=(-\infty;3]\setminus[2;5)=(-\infty;2)$. Các số nguyên không âm thuộc tập  này là $0$ và $1$. Có 2 số nguyên không âm.
			\itemch $B\setminus A=[2;5)\setminus(-\infty;3]=(3;5)$. $C=\{x|x\in\mathbb{N}, x\in(3;5)\}=\{4\}$. Tập C có  1 phần tử nên có $2^{1}=2$ tập con ($\emptyset$ và $\{4\}$).
		\end{itemchoice}
	 }
\end{ex}

% nguồn: MaHoaNH_DungSai_Lop10
\begin{ex}%[0D1H3-4]%[ma:bc823e38]
	Cho tam giác $ABC$ có $AB=4$, $AC=3$ và $\widehat{BAC}=60^{\circ}$. (Kết quả làm  tròn đến hàng đơn vị).
	\choiceTF {\True  $BC=\sqrt{13}$}{ $\widehat{ABC}\approx75^{\circ}$}{ $\widehat{ACB}\approx45^{\circ}$}{\True $AB$ là cạnh có độ dài lớn nhất trong tam giác}
	\loigiai{ 
		\begin{itemchoice}
			\itemch Áp  dụng định lí Côsin: $BC^{2}=4^{2}+3^{2}-2\cdot4\cdot3\cdot\cos 60^{\circ}=16+9-12=13$. Do đó $BC=\sqrt{13}$.
			\itemch Áp dụng định lí Côsin: $\cos B=\dfrac{AB^{2}+BC^{2}-AC^{2}}{2AB\cdot BC}=\dfrac{4^{2}+13-3^{2}}{2\cdot4\cdot\sqrt{13}}=\dfrac{20}{8\sqrt{13}}=\dfrac{5\sqrt{13}}{26}$. Do  đó $\widehat{ABC}\approx46^{\circ}$.
			 \itemch $\widehat{ACB}=180^\circ-\widehat{BAC}-\widehat{ABC}\approx180^\circ-60^\circ-46^\circ\approx74^\circ$.
			\itemch Ta có $AB=4$, $AC=3$, $BC=\sqrt{13}\approx3{,}6$. Cạnh lớn nhất là  $AB=4$.
		\end{itemchoice}
	} 
\end{ex}

% nguồn: MaHoaNH_DungSai_Lop10
\begin{ex}%[0D1H3-4]%[ma:1465c22a]
	Cho hai tập hợp $A=\{x\in  \mathbb{R}\mid -5\le x\le-3\}$, $B=\left(-3;2\right)$. Khi đó
	\choiceTF 
		 {\True $A=\left[-5;-3\right]$}
		{\True $A\cup  B=\left[-5;2\right)$}
		 {$A\cap B=\{-3\}$}
		{\True $\mathrm{C}_{\mathbb{R}} \left(A\cup B\right)=\left(-\infty;-5\right)\cup  \left[2;+\infty \right)$}
	 \loigiai{
		\begin{itemchoice}
			\itemch $A=\{x\in \mathbb{R}\mid -5\le x\le-3\}\Rightarrow A=\left[-5;-3\right]$. 
			 \itemch $A\cup B=\left[-5;2\right)$.
			\itemch $A\cap  B=\varnothing$.
			\itemch $\mathrm{C}_{\mathbb{R}} \left(A\cup B\right)=\mathbb{R}\setminus \left[-5;2\right)=\left(-\infty;-5\right)\cup \left[2;+\infty  \right)$.
		\end{itemchoice}
	} 
\end{ex}

% nguồn: MaHoaNH_DungSai_Lop10
\begin{ex}%[0D1H3-4]%[ma:415441b4]
	Cho hai tập hợp $A=\{x \in \mathbb{R} \mid-1<  x \leq 4\}$ và $B=(-3;+\infty)$.
	 Xác đinh tính đúng, sai các khẳng đinh sau
	 \choiceTF
		{Tâp hợp $A$  có $5$ phần từ}
		{$A\cap B=(-3; 4]$} 
		 {\True $A\setminus B=\varnothing$}
		 {\True $C_BA=(-3;-1] \cup(4;+\infty)$}
	 \loigiai{
		\begin{itemchoice}
			\itemch Sai: $A=\{x \in \mathbb{R} \mid-1< x \leq 4\} \Rightarrow A=(-1; 4]$  nên tập $A$ có vô số phần tử.
			\itemch Sai: $A=(-1; 4]$ và $B=(-3;+\infty)  \Rightarrow A\subset B\Rightarrow A\cap B=A=(-1; 4]$.
			\itemch Đúng: vì $A\subset  B\Rightarrow A\setminus B=\varnothing$.
			\itemch Đúng: ta có  $A\subset B\Rightarrow C_BA=(-3;-1] \cup(4;+\infty)$.
		\end{itemchoice}
	} 
\end{ex}


%%%%%%%%%%%%%%%% TRẢ LỜI NGẮN %%%%%%%%%%%%%%%%

% nguồn: MaHoaNH_TraLoiNgan_Lop10
\begin{bt}%[0D1V3-4]%[ma:bd954b33]
	Cho các tập hợp $A=\left(2;+\infty \right)$ và $B=\left[m^2-7;+\infty \right)$ với  $m > 0$. Số giá trị nguyên của $m$ để $A\setminus B$ là một khoảng $\left(a; b\right)$ thoả mãn $b-a$ thuộc đoạn $\left[3; 16\right]$ là?
	\shortans{2} 
	 \loigiai{
		Điều kiện để $A\setminus  B\ne \varnothing$ là $\heva{&m^2-7> 2\\
		&m > 0} \Leftrightarrow  \heva{&m^2> 9\\
		&m > 0} 
		\Leftrightarrow m  > 3$.\\
		 Khi đó $A\setminus B=\left(2;m^2-7\right)\Rightarrow a=2$, $b=m^2-7$.\\
		Ta có $b-a$ thuộc đoạn $\left[3; 16\right]\Leftrightarrow 3\le m^2-9\le 16\Leftrightarrow 12\le  m^2\le 25$.\\
		Do $m > 3$ và $m\in \mathbb{Z}$ nên $m\in \left\{4; 5\right\}$.\\ Vậy có  $2$ giá trị $m$ nguyên thoả mãn.
	} 
\end{bt}


%%%%%%%%%%%%%%%% TỰ LUẬN %%%%%%%%%%%%%%%%

% nguồn: MaHoaNH_TuLuan_Lop10
\begin{bt}%[0D1N3-4]%[ma:481226fd]
	Cho hai tập hợp $A=\{n \in \mathbb{N} \mid-2<n \leqslant 4\}$ và $B=\{-1; 0\}$.  Hãy xác định tập hợp $A \cap B$, $A \setminus B$.
	 \loigiai{
		Ta có $A=\{0; 1; 2; 3; 4\}$,  $B=\{-1; 0\}$.\\
		Suy ra $A  \cap B=\{0\}$ và $A \setminus B=\{1; 2; 3; 4\}$.
	 }
\end{bt}

% nguồn: MaHoaNH_TuLuan_Lop10
\begin{bt}%[0D1N3-4]%[ma:18fd1a47]
	Cho hai tập hợp $A=\{n \in \mathbb{Z} \mid -2 < n \leq 4\}$ và $B=\{-1; 0;1;5\}$. Hãy xác định  tập hợp $A \cup B$, $A \setminus B$.
	\loigiai{ 
		Ta có $A=\{-1; 0; 1; 2; 3; 4\}$, $B=\{-1;  0;1;5\}$.\\
		Suy ra $A \cup  B=\{-1; 0; 1; 2; 3; 4; 5\}$ và $A \setminus B=\{2; 3; 4\}$.
	} 
\end{bt}

% nguồn: MaHoaNH_TuLuan_Lop10
\begin{bt}%[0D1N3-4]%[ma:e1a629c7]
	 Cho các tập hợp $A=\{0;1;2;3;4;5;6\}$, $B=\{-3;-1;1;2;3\}$ và $C=\{x\in \mathbb{N}\mid x \,\text{là ước của}\, 6\}$.
	\begin{listEX}
		\item  Xác định tập $A\cup B$.
		\item Tập $C\setminus  B$ có bao nhiêu phần tử?
	\end{listEX}
	\loigiai{Ta có  $A=\{0;1;2;3;4;5;6\}$, $B=\{-3;-1;1;2;3\}$, $C=\{1;2;3;6\}$.
		\begin{listEX}
			\item  Ta có $A\cup B=\{-3;-1;0;1;2;3;4;5;6\}$.
			\item  Ta có $C\setminus B=\{2;3\}$.\\
			Vậy $C\setminus  B$ có $2$ phần tử.
		\end{listEX}
	} 
\end{bt}

% nguồn: MaHoaNH_TuLuan_Lop10
\begin{bt}%[0D1H3-4]%[ma:0b7d78dc]
	Cho tập hợp $A$ và $B$ thỏa $C_{\mathbb{R}} A=\left[-10; \sqrt{2019}\right)$ và $C_{\mathbb{R}} B=\left(-12; \sqrt{2018}\right]$. Xác định tập hợp $C_{\mathbb{R}}(A \cap B)$. 
	\loigiai{ 
		Do $C_{\mathbb{R}}  A=\left[-10; \sqrt{2019}\right)$, suy ra $A=(-\infty;-10) \cup \left[\sqrt{2019};+\infty\right)$.\\
		Do $C_{\mathbb{R}} B=\left(-12; \sqrt{2018}\right]$,  suy ra $B=(-\infty;-12] \cup \left(\sqrt{2018};+\infty\right]$.\\
		Suy ra  $A \cap B=(-\infty;-12] \cup \left[\sqrt{2019};+\infty\right)$.\\
		Vậy $C_{\mathbb{R}}(A  \cap B)=\left(-12; \sqrt{2019}\right)$.
	} 
\end{bt}

% nguồn: MaHoaNH_TuLuan_Lop10
\begin{bt}%[0D1H3-4]%[ma:b3f10ff6]
	Cho các tập hợp $A=(2;+\infty)$ và $B=\left[m^2-7;+\infty\right)$  với $m>0$. Tìm tất cả các số thực $m$ để $A \setminus B$ là một khoảng có độ dài bằng $16$.
	\loigiai{ 
		 Điều kiện để $A \setminus B \neq \varnothing$ là $\heva{
		&m^2-7>2  \\
		&m>0}  \Leftrightarrow \heva{
		&m^2>9  \\
		&m>0} \Leftrightarrow m>3$.\\ 
		Khi đó $A  \setminus B=\left(2; m^2-7\right)$.\\
		Độ dài khoảng $A \setminus B$ bằng $16 \Leftrightarrow m^2-7-2=16 \Rightarrow  m=5$ (do $m>3$).
	} 
\end{bt}

% nguồn: MaHoaNH_TuLuan_Lop10
\begin{bt}%[0D1H3-4]%[ma:3ee08ede]
	Cho khoảng  $A=\left(-\infty; \dfrac{6}{2-m}\right)$ và khoảng $B=(1-m;+\infty)$. Tìm tất cả các giá trị thực của tham số $m$ để $A \setminus B=A$.
	\loigiai{ 
		Ta có 
		\begin{eqnarray*}
			A \setminus  B=A \Leftrightarrow A \cap B=\varnothing &\Leftrightarrow& \dfrac{6}{2-m} \leq 1-m\\
			&\Leftrightarrow& \dfrac{-m^{2}+3 m+4}{2-m}  \leq 0\\
			&\Leftrightarrow&  \hoac{
			&m \leq-1 \\ 
			&2<m  \leq 4.}
		\end{eqnarray*}
	 }
\end{bt}

% nguồn: MaHoaNH_TuLuan_Lop10
\begin{bt}%[0D1H3-4]%[ma:c6264326]
	Cho các tập hợp$A=\{x \in \mathbb{R} \mid x<3\}$, $B=\{x \in \mathbb{R} \mid 1<x \leq 5\}$ và $C=\{x \in \mathbb{R} \mid-2 \leq  x \leq 4\}$.
	\begin{listEX}
		\item Hãy viết lại các tập hợp $A$,  $B$, $C$ dưới kí hiệu khoảng, nửa khoảng, đoạn.
		\item Tìm $A \cup B$, $A \cap B$, $A  \setminus B$.
		\item Tìm $C_{\mathbb{R}} A$, $C_{\mathbb{R}} B$, $C_{\mathbb{R}}  C$.
		\item Tìm  $(B \cup C) \setminus (A \cap C)$.
	\end{listEX}
	\loigiai{ 
		\begin{listEX}
			\item  Ta có $A=(-\infty; 3)$, $B=(1; 5]$, $C=[-2; 4]$.
			\item Suy ra  $A \cup B=(-\infty; 5]$.\\
			Suy ra $A \cap B=(1; 3)$.\\ 
			Suy ra $A \setminus B=(-\infty; 1]$.\\ 
			\item Ta có $C_{\mathbb{R}} A=\mathbb{R} \setminus A=[3;+\infty)$, $C_{\mathbb{R}} B=\mathbb{R} \setminus B=(-\infty; 1]  \cup(5;+\infty)$, \\$C_{\mathbb{R}} C=\mathbb{R} \setminus C=(-\infty;-2) \cup(4;+\infty)$.
			\item Ta có  $A \cap C=[-2; 3)$ và $B \cup C=[-2; 5]$.\\
			 Suy ra $(B \cup C) \setminus (A \cap C)=[3; 5]$.
		\end{listEX}
	 }
\end{bt}

% nguồn: MaHoaNH_TuLuan_Lop10
\begin{bt}%[0D1H3-4]%[ma:f2b1bc48]
	Cho tập hợp $A=\{x\in \mathbb{R}\mid-3<x<3\},B=\{x\in  \mathbb{R}\mid-2<x\leq 3\}$ và $C=\{x\in \mathbb{R}\mid 0\leq x\leq 4\}$. Hãy tìm tập hợp $D$ thỏa mãn
	\begin{enumEX}{3}
		\item  $D=(A\cup B)\cup C$;
		\item $D=(A\cup B)\cap  C$;
		\item $D=(A\cap B)\cap  C$;
		\item $D=(A\cap B)\cup C$; 
		\item  $D=(A\cap B)\backslash C$;
		\item  $D=(B\cup C)\backslash A$;
		\item $D=(B\backslash A)\cup(C\backslash A)$; 
		 \item $D=(B\backslash A)\backslash C$;
		\item $D=(B\backslash A)\cup C$; 
		\item $D=(A\backslash B)\cup(A\backslash C)$. 
	\end{enumEX}
	\loigiai{ 
		\begin{enumerate}
			\item Ta  có $A\cup B=\{x\in \mathbb{R}\mid-3<x\leq 3\}$; $D=(A\cup B)\cup C=\{x\in \mathbb{R}\mid-3<x\leq 4\}$.
			\item  Ta có $A\cup B=\{x\in \mathbb{R}\mid-3<x\leq 3\}$; $D=(A\cup B)\cap C=\{x\in \mathbb{R}\mid 0<x\leq 3\}$.
			 \item Ta có $A\cap B=\{x\in \mathbb{R}\mid-2<x<3\}$; $D=(A\cap B)\cap C=\{x\in \mathbb{R}\mid 0<x<3\}$.
			\item Ta có $A\cap B=\{x\in \mathbb{R}\mid-2<x<3\}$; $D=(A\cap B)\cup C=\{x\in  \mathbb{R}\mid-2<x\leq 4\}$.
			\item Ta có $A\cap B=\{x\in \mathbb{R}\mid-2<x<3\}$; $D=(A\cap B)\backslash C=\{x\in \mathbb{R}\mid-2<x\leq 0\}$. 
			\item Ta có  $B\cap C=\{x\in \mathbb{R}\mid-2<x\leq 4\}$ nên $D=(B\cup C)\backslash A=\{x\in \mathbb{R}\mid 3\leq x\leq 4\}$.
			\item Ta có $B\backslash A=\{3\}$; $C\backslash A=\{x\in \mathbb{R}\mid 3\leq x\leq 4\}$ nên khi đó ta có\[D=(B\backslash A)\cup(C\backslash A)=C\backslash A=\{x\in  \mathbb{R}\mid 3\leq x\leq 4\}.\]
			\item Ta có $B\backslash A=\{3\}$ nên $D=(B\backslash A)\backslash  C=\varnothing$.
			\item Theo h)  thì $D=(B\backslash A)\cup C=\{x\in \mathbb{R}\mid 0\leq x\leq 4\}$.
			 \item Ta có $A\backslash B=\{x\in \mathbb{R}\mid-3<x\leq-2\}$; $A\backslash C=\{x\in \mathbb{R}\mid-3<x\leq 0\}$.
		\end{enumerate}
	} 
\end{bt}
